% Page 026

\seite{26}

\abschnitt{Kaisergeburtstag}
\pstart
Die Kaisergeburtstagsfeier wurde gemeinschaftlich mit\ob
Schülern der Pfarrei in Seffern abgehalten. Es entspricht\ob
einer alten Gepflogenheit. Die Festrede hielt\ob
Herr Lehrer Koch aus Heilenbach. Der\ob
Vortrag drehte sich um die drei Punkte: für\ob
unsern Landesfürsten\ob
1.\ lieben,\ob
2.\ ehren\ob
3.\ ihm gehorsam sein.

Dem Ortsschulinspektor Herrn Pfarrer Eltges wurde\ob
ein Hoch ausgebracht, worauf die Verteilung der Kaiserwecken\ob
folgte.

\pend
\abschnitt{Osterprüfung}
\pstart

Die Osterprüfung wurde am 29.\ März durch den Orts\ohb
schulinspektor abgehalten.

Am 4.\ April wurde die Schule vom Herrn Kreisschul\ohb
inspektor revidiert. Dr. Reuter visitirt.

Mit dem Schlusse des Schuljahres am 6.\ April\ob
wurden fünf Schüler der Schule entlassen. Sämt\ohb
liche Schüler traten in das Geschäft bzw.\ in den\ob
Beruf des Vaters ein.

Mit Beginn des Schuljahres traten zwei Söhne hiesiger\ob
Bürger in das Gymnasium zu Prüm ein, um sich\ob
zum Priesterstande vorzubereiten. Privatim vorgebildet\ob
wurden sie auf der 4.\ Klasse (Untertertia) auf\ohb
genommen.
\pend
